% Google Student Researcher — tailored 1-page resume
% Compile from resumes/google-student-researcher/: pdflatex Rojae_Mighty_Resume.tex
\documentclass[11pt,letterpaper]{article}

\usepackage[margin=0.6in]{geometry}
\usepackage{enumitem}
\usepackage[hidelinks]{hyperref}
\usepackage{titlesec}
\usepackage{tabularx}

\pagestyle{empty}
\setlength{\parindent}{0pt}
\setlength{\parskip}{0pt}
\pagenumbering{gobble}

\titleformat{\section}{\large\bfseries}{}{0pt}{}[\titlerule]
\titlespacing*{\section}{0pt}{8pt}{4pt}

\newcommand{\header}[1]{{\LARGE\bfseries #1}}
\newcommand{\lineitem}[2]{%
  \textbf{#1} \hfill \textit{#2}\\[-2pt]
}
\newcommand{\subheader}[1]{\textit{#1}\\[2pt]}
\newcommand{\cvitem}[1]{\item #1}
\setlist[itemize]{leftmargin=*, topsep=4pt, itemsep=5pt, parsep=0pt}

\begin{document}

\header{Rojae Mighty}\\[2pt]
\href{mailto:rojae.mighty@stonybrook.edu}{rojae.mighty@stonybrook.edu} \quad$\cdot$\quad
Stony Brook, NY \quad$\cdot$\quad
\href{https://rojaemighty.com}{rojaemighty.com} \quad$\cdot$\quad
\href{https://www.linkedin.com/in/rojaemight/}{linkedin.com/in/rojaemight}

\section{Education}
\lineitem{Stony Brook University}{Expected Graduation: 5/28}
B.S.\ in Physics (Natural Sciences)\\
Programming languages: \textbf{Python}, \textbf{C/C++} \quad$\cdot$\quad
Relevant coursework: statistical mechanics, computational methods, data analysis, probability

\section{Research Experience}

\lineitem{Brookhaven National Laboratory}{2025--Present}
\begin{itemize}
  \cvitem{\textbf{Department of Energy Science Undergraduate Laboratory Internship (SULI)} (Summers 2025, 2026): Electron-ion collision studies of semi-inclusive processes; implemented a spatially dependent parton distribution function using \textbf{ROOT}, \textbf{Python}, and \textbf{Unix/Linux shell scripting}.}
\end{itemize}

\lineitem{Center for Frontiers in Nuclear Science (CFNS)}{2025--Present}
\subheader{Undergraduate Researcher \quad Advisor: Zhoudunming (Kong) Tu}
\begin{itemize}
  \cvitem{Developed \textbf{machine learning} pipelines in \textbf{Python} with \textbf{scikit-learn} for kinematic reconstruction on large-scale collision datasets; applied \textbf{data science} methods including \textbf{feature engineering} and \textbf{statistical data analysis}.}
  \cvitem{Used \textbf{C++} and \textbf{ROOT} with \textbf{Monte Carlo event generators} to study gluon saturation through geometric scaling; ran \textbf{numerical simulation} studies of high-energy nuclear collision phenomena.}
  \cvitem{Delivered \textbf{conference presentations} at DIS 2026, APS Global Physics Summit (oral), and Stony Brook Physics Colloquium; prepared materials in \textbf{LaTeX}.}
\end{itemize}

\section{Projects}

\lineitem{Market Microstructure Forecasting}{2025--Present \quad \href{https://github.com/iRojae/market-microstructure-forecasting}{GitHub}}
\begin{itemize}
  \cvitem{Built a \textbf{Python} \textbf{machine learning} pipeline with \textbf{pandas}, \textbf{NumPy}, and \textbf{scikit-learn}; \textbf{feature engineering} and \textbf{walk-forward validation} on time-series limit order book data; \textbf{Git}-versioned on GitHub.}
\end{itemize}

\lineitem{J.P. Morgan Quantitative Research (\href{https://www.theforage.com/simulations/jpmorgan/quantitative-research-11oc}{Forage})}{2026}
\begin{itemize}
  \cvitem{Forage certificate: commodity \textbf{statistical data analysis}; \textbf{derivatives pricing}; \textbf{Python} credit-risk model; FICO bucketing via \textbf{dynamic programming}/\textbf{scikit-learn} (MSE/log-likelihood).}
\end{itemize}

\section{Honors \& Presentations}
\begin{itemize}
  \cvitem{1st Place, Physical Sciences Presentation, Yale Undergraduate Research Conference (2026)}
  \cvitem{Outstanding Oral Presentation Award, National Society of Black Physicists Conference (2025)}
  \cvitem{Barry Barish Undergraduate Research Travel Award, Stony Brook University (2026)}
  \cvitem{\textit{Geometric Scaling and Gluon Saturation at the EIC} --- DIS 2026, Yale URC, NSBP, APS DNP (CEU), APS Global Physics Summit, Stony Brook Colloquium}
\end{itemize}

\section{Technical Skills}
\begin{tabularx}{\textwidth}{@{}>{\bfseries}l X@{}}
Languages: & Python, C/C++, Unix/Linux shell scripting \\
Libraries: & pandas, NumPy, scikit-learn \\
ML \& data: & machine learning, data science, statistical data analysis, feature engineering, walk-forward validation, numerical simulation, dynamic programming \\
Tools: & ROOT, Git, LaTeX, Monte Carlo event generators
\end{tabularx}

\end{document}
